\documentclass[12pt]{article}
\usepackage[T2A]{fontenc}
\usepackage[utf8]{inputenc}
\usepackage[russian]{babel}
\usepackage{amsmath}
\usepackage{amsfonts}
\usepackage{amssymb}
\usepackage{array}
\usepackage{geometry}
\geometry{margin=2.2cm}

\title{WonderFullyHE: математическое описание API}
\date{}

\begin{document}
\maketitle

\section*{Назначение библиотеки}

WonderFullyHE реализует библиотечный слой для защищённых вычислений
над вещественными данными на базе схемы CKKS из Microsoft SEAL.
Публичный API скрывает низкоуровневые типы SEAL и предоставляет:
\begin{itemize}
    \item настройку CKKS-профилей;
    \item кодирование, шифрование, вычисления, расшифрование и декодирование;
    \item контроль численной ошибки;
    \item ABFT-проверки корректности вычислений;
    \item rotation-based линейные преобразования;
    \item вычисление полиномов от ciphertext;
    \item диагностику вычислительной глубины и bootstrapping-конвейер.
\end{itemize}

\section*{CKKS-профиль}

Профиль CKKS задаётся набором
\[
    P = (N, Q, \Delta, S),
\]
где:
\begin{itemize}
    \item $N$ --- степень полиномиального модуля;
    \item $Q = (q_0, q_1, \ldots, q_l)$ --- цепочка коэффициентных модулей,
    в коде задаётся битовыми размерами;
    \item $\Delta$ --- масштаб CKKS;
    \item $S$ --- число используемых слотов, $S \leq N/2$.
\end{itemize}

Публичный модуль \texttt{m2424::profiles} задаёт готовые профили,
чтобы пользователь не собирал параметры CKKS вручную:

\[
\begin{array}{|c|c|c|c|c|c|}
\hline
\text{Профиль} & N & S & Q\text{, бит} & \Delta & \text{Оценка глубины} \\
\hline
\texttt{fast\_demo\_ckks} & 4096 & 2048 & 40-30-30 & 2^{30} & 1 \\
\hline
\texttt{basic\_ckks} & 8192 & 4096 & 60-40-40-60 & 2^{40} & 2 \\
\hline
\texttt{balanced\_ckks} & 8192 & 4096 & 50-40-40-40-48 & 2^{40} & 3 \\
\hline
\texttt{depth\_ckks} & 16384 & 8192 & 60-40-40-40-40-60 & 2^{40} & 4 \\
\hline
\texttt{high\_precision\_ckks} & 16384 & 8192 & 60-50-50-50-60 & 2^{50} & 3 \\
\hline
\end{array}
\]

\texttt{fast\_demo\_ckks} нужен для быстрых локальных запусков.
\texttt{basic\_ckks} используется в основных демонстрациях.
\texttt{balanced\_ckks} даёт дополнительный уровень умножения при
той же степени \(N=8192\). \texttt{depth\_ckks} используется для
анализа глубины и bootstrapping-блоков. \texttt{high\_precision\_ckks}
увеличивает масштаб до \(2^{50}\), когда приоритетом является меньшая
ошибка округления, а не скорость.

Готовые профили являются пресетами. Целевой способ выбора параметров ---
\texttt{CkksParameterPlanner}:
\[
    (\varepsilon, D, S, \mathsf{ops}, \mathsf{security})
    \mapsto
    (N, Q, \Delta),
\]
где \(\varepsilon\) --- целевая ошибка, \(D\) --- требуемая
мультипликативная глубина, \(S\) --- число слотов.

Для целевой ошибки используется оценка:
\[
    b_{\mathsf{result}} = \lceil -\log_2(\varepsilon) \rceil,
\]
\[
    b_{\mathsf{work}}
    =
    b_{\mathsf{result}}
    +
    b_{\mathsf{loss}}(\mathsf{ops}),
\]
где \(b_{\mathsf{loss}}\) задаётся калибровкой конкретного набора
операций. Для \(\varepsilon=10^{-9}\) имеем
\[
    b_{\mathsf{result}} = 30.
\]
Текущий benchmark для двух последовательных
\texttt{mul\_relin\_rescale} даёт потерю порядка 14 бит, поэтому
практический быстрый режим выбирает рабочие модули около 45 бит и
\(\log_2(\Delta)\approx45\).

\section*{Адаптер над Microsoft SEAL}

Создание контекста:
\[
    \texttt{SealAdapter::create}(P) \mapsto \mathsf{SEALContext}(P).
\]

Генерация ключей:
\[
    \texttt{keygen}(r,g) \mapsto
    (\mathsf{sk}, \mathsf{pk}, \mathsf{rlk}^{(r)}, \mathsf{gk}^{(g)}),
\]
где $r,g \in \{0,1\}$ задают необходимость Relin- и Galois-ключей.

Для оптимизации размера Galois-ключей поддерживается генерация только
заданного набора ротаций:
\[
    \texttt{keygen}(K,r) \mapsto
    (\mathsf{sk}, \mathsf{pk}, \mathsf{rlk}^{(r)}, \{\mathsf{gk}_k\}_{k \in K}),
\]
где \(K\) --- множество rotation steps.

Сериализация объектов:
\[
    \texttt{save\_public\_key}() \mapsto B_{\mathsf{pk}},
    \quad
    \texttt{save\_secret\_key}() \mapsto B_{\mathsf{sk}},
\]
\[
    \texttt{save\_relin\_keys}() \mapsto B_{\mathsf{rlk}},
    \quad
    \texttt{save\_galois\_keys}() \mapsto B_{\mathsf{gk}},
    \quad
    \texttt{save\_cipher}(c) \mapsto B_c.
\]
Обратные операции \texttt{load\_*} восстанавливают объект в новом
CKKS-контексте с тем же профилем. Это позволяет разделять контексты:
контекст с секретным ключом выполняет расшифрование, а вычислительный
контекст получает только ciphertext и evaluation keys.

Кодирование и шифрование:
\[
    \texttt{encode}(\mathbf{x}) =
    \mathsf{Encode}_{\Delta}(\mathbf{x}) = p,
\]
\[
    \texttt{encrypt}(p) =
    \mathsf{Enc}_{\mathsf{pk}}(p) = c.
\]

Расшифрование и декодирование:
\[
    \texttt{decrypt}(c) =
    \mathsf{Dec}_{\mathsf{sk}}(c) = p,
\]
\[
    \texttt{decode}(p) =
    \mathsf{Decode}_{\Delta}(p) = \hat{\mathbf{x}}.
\]

Для операций ciphertext--plaintext используется кодирование на уровне
заданного ciphertext:
\[
    \texttt{encode\_like}(\mathbf{a},c)
    =
    \mathsf{Encode}_{\mathsf{scale}(c),\mathsf{level}(c)}(\mathbf{a}).
\]

\section*{Гомоморфные операции}

Сложение:
\[
    \texttt{add}(c_1,c_2) = c_1 + c_2.
\]

Вычитание:
\[
    \texttt{sub}(c_1,c_2) = c_1 - c_2.
\]

Сложение и вычитание с plaintext:
\[
    \texttt{add\_plain}(c,p) = c + p,
    \qquad
    \texttt{sub\_plain}(c,p) = c - p.
\]

Умножение ciphertext на plaintext с rescale:
\[
    \texttt{mul\_plain\_rescale}(c,p)
    =
    \mathsf{Rescale}(c \cdot p).
\]

Умножение с релинеаризацией и rescale:
\[
    \texttt{mul\_relin\_rescale}(c_1,c_2)
    =
    \mathsf{Rescale}
    \left(
        \mathsf{Relin}_{\mathsf{rlk}}(c_1 \cdot c_2)
    \right).
\]

При умножении масштаб растёт:
\[
    \Delta \cdot \Delta = 2^{40} \cdot 2^{40} = 2^{80}.
\]

Операция \(\mathsf{Rescale}\) делит результат на следующий
примерно 40-битный модуль и возвращает масштаб к значению, близкому к
\(2^{40}\). При этом расходуется один уровень modulus chain.

Ротация CKKS-слотов:
\[
    \texttt{rotate}(c,k) = \mathsf{Rot}_k(c).
\]

Для сложения результатов разных степеней поддерживается выравнивание
уровня и масштаба:
\[
    \texttt{match\_level\_and\_scale}(c,t)
    \mapsto c^*:
    \mathsf{level}(c^*)=\mathsf{level}(t),
    \quad
    \mathsf{scale}(c^*)=\mathsf{scale}(t).
\]

\section*{Линейные преобразования через ротации}

Модуль \texttt{LinearTransform} применяет преобразование:
\[
    L(c)
    =
    \sum_{i=0}^{m-1}
    a_i \cdot \mathsf{Rot}_{k_i}(c),
\]
где \(a_i\) --- plaintext-коэффициенты, а \(k_i\) --- шаги ротации.

В коде каждый член вычисляется как:
\[
    t_i =
    \texttt{mul\_plain\_rescale}
    \left(
        \texttt{rotate}(c,k_i),
        \texttt{encode\_like}(a_i,\texttt{rotate}(c,k_i))
    \right),
\]
после чего члены складываются на одном уровне.

Функция \texttt{sum\_slots} считает сумму первых \(m\) слотов:
\[
    s =
    \sum_{i=0}^{m-1} x_i.
\]
Для \(m=2^d\) используется последовательность ротаций:
\[
    1,2,4,\ldots,2^{d-1}.
\]
Результат суммы помещается в первый CKKS-слот.

\section*{Вычисление полиномов}

Модуль \texttt{PolynomialEvaluator} вычисляет полином:
\[
    P(c)
    =
    \sum_{j \in D} \alpha_j c^j,
\]
где \(D\) --- набор используемых степеней, а \(\alpha_j\) ---
вещественные коэффициенты.

Степени ciphertext строятся последовательными умножениями:
\[
    c^{j+1}
    =
    \texttt{mul\_relin\_rescale}
    \left(
        c^j,
        \texttt{match\_level\_and\_scale}(c,c^j)
    \right).
\]

Каждый член полинома умножается на plaintext-коэффициент:
\[
    \alpha_j c^j
    =
    \texttt{mul\_plain\_rescale}
    \left(
        c^j,
        \texttt{encode\_scalar\_like}(\alpha_j,c^j)
    \right).
\]

Перед сложением членов выполняется выравнивание уровня и масштаба.
Этот блок используется как программная основа для этапа \(\mathsf{EvalMod}\).

\section*{Диагностика ciphertext}

Для ciphertext \(c\) библиотека возвращает набор параметров:
\[
    \texttt{info}(c) =
    (\mathsf{scale}(c), \mathsf{chain\_index}(c),
    \mathsf{coeff\_modulus\_size}(c), \mathsf{ciphertext\_size}(c)).
\]

Также измеряется сериализованный размер:
\[
    \texttt{serialized\_size}(c) = |\mathsf{Serialize}(c)|.
\]

Эти значения используются в benchmark, анализе глубины и bootstrapping-отчёте.

\section*{Контроль точности}

Пусть \(\mathbf{y}\) --- эталонный результат CPU-вычисления,
а \(\hat{\mathbf{y}}\) --- результат после гомоморфного вычисления,
расшифрования и декодирования. Для векторов одинаковой длины \(n\):

\[
    \mathsf{max\_abs\_error}
    =
    \max_{0 \leq i < n}|y_i - \hat{y}_i|,
\]

\[
    \mathsf{mean\_abs\_error}
    =
    \frac{1}{n}\sum_{i=0}^{n-1}|y_i - \hat{y}_i|.
\]

Критерий прохождения проверки:
\[
    \mathsf{ok}
    =
    \mathsf{max\_abs\_error} \leq \mathsf{tolerance}.
\]

Для доказуемого контракта точности ошибка должна распространяться
по стадиям:
\[
    E_{\mathsf{total}}
    \leq
    E_{\mathsf{encode}}
    +
    E_{\mathsf{encrypt}}
    +
    E_{\mathsf{rescale}}
    +
    E_{\mathsf{keyswitch}}
    +
    E_{\mathsf{poly}}
    +
    E_{\mathsf{linear}}.
\]
Планировщик должен запускать вычисление только если рассчитанная
верхняя граница не превышает заданную \(\mathsf{tolerance}\).

\section*{ABFT-проверки}

Контрольная сумма полезной нагрузки:
\[
    \mathsf{checksum}(\mathbf{x})
    =
    \sum_{i=0}^{m-1}x_i.
\]

Добавление checksum-слота:
\[
    \mathbf{x}' =
    (x_0,x_1,\ldots,x_{m-1},\mathsf{checksum}(\mathbf{x})).
\]

Проверка добавленного checksum-слота:
\[
    \mathsf{expected} =
    \sum_{i=0}^{m-1} y_i,
    \qquad
    \mathsf{observed} = y_m,
\]
\[
    \mathsf{abs\_error}
    =
    |\mathsf{expected} - \mathsf{observed}|,
    \qquad
    \mathsf{ok}
    =
    \mathsf{abs\_error} \leq \mathsf{tolerance}.
\]

Для сложения и вычитания используются линейные свойства:
\[
    \mathsf{checksum}(\mathbf{x}+\mathbf{y})
    =
    \mathsf{checksum}(\mathbf{x})+
    \mathsf{checksum}(\mathbf{y}),
\]
\[
    \mathsf{checksum}(\mathbf{x}-\mathbf{y})
    =
    \mathsf{checksum}(\mathbf{x})-
    \mathsf{checksum}(\mathbf{y}).
\]

Для ротации используется инвариант суммы:
\[
    \mathsf{checksum}(\mathsf{Rot}_k(\mathbf{x}))
    =
    \mathsf{checksum}(\mathbf{x}).
\]

Для умножения текущая проверка сравнивает сумму результата
с CPU-эталоном покомпонентного произведения:
\[
    \mathsf{expected}
    =
    \sum_{i=0}^{m-1} x_i y_i,
    \qquad
    \mathsf{observed}
    =
    \sum_{i=0}^{m-1} \widehat{(xy)}_i.
\]

\section*{Анализ вычислительной глубины}

В профиле \texttt{depth\_ckks} библиотека анализирует цепочку
последовательных зашифрованных возведений в квадрат:
\[
    x \rightarrow x^2 \rightarrow x^4 \rightarrow x^8 \rightarrow x^{16}.
\]

Каждый шаг вызывает:
\[
    c_{i+1} =
    \texttt{mul\_relin\_rescale}(c_i,c_i).
\]

После каждого успешного шага фиксируются:
\[
    \mathsf{scale}(c_i),\quad
    \mathsf{chain\_index}(c_i),\quad
    \mathsf{serialized\_size}(c_i),\quad
    \mathsf{error}(c_i).
\]

На текущем профиле следующий переход к \(x^{32}\) останавливается,
потому что доступные уровни modulus chain для очередного rescale
исчерпаны.

\section*{Bootstrapping-модуль}

Bootstrapping-модуль выделен как отдельный компонент библиотеки.
Он связан с анализом глубины и фиксирует конвейер:
\[
    \mathsf{ModRaise}
    \rightarrow
    \mathsf{CoeffToSlot}
    \rightarrow
    \mathsf{EvalMod}
    \rightarrow
    \mathsf{SlotToCoeff}.
\]

Текущий scalable target --- не плотная таблица преобразования, а
factorized FFT-like реализация \(\mathsf{CoeffToSlot}\) и
\(\mathsf{SlotToCoeff}\). Плотное диагональное разложение остаётся
reference-path для малых размеров.

Для \(\mathsf{EvalMod}\) при нормализации
\[
    |u| \leq 2^{-10}
\]
достаточен полином \(P3\):
\[
    P3(u) = u - 6.579736267393u^3.
\]
Его ошибка аппроксимации на указанном интервале существенно меньше
\(10^{-9}\); более высокие степени \(P5/P7\) увеличивают глубину и
используются только при расширении интервала или более жёсткой цели.

Цель итогового преобразования:
\[
    \mathsf{Bootstrap}(c) = c',
\]
где должны выполняться условия:
\[
    \mathsf{Dec}(c') \approx \mathsf{Dec}(c),
\]
\[
    \mathsf{level}(c') > \mathsf{level}(c).
\]

В публичном отчёте \texttt{BootstrapReport} фиксируются:
\begin{itemize}
    \item параметры входного ciphertext;
    \item параметры ciphertext на границе вычислительной глубины;
    \item число успешно выполненных умножений;
    \item следующий показатель степени, на котором вычисление останавливается;
    \item причина остановки;
    \item стадии конвейера и их статусы;
    \item состояние критериев \(\mathsf{Dec}(c') \approx \mathsf{Dec}(c)\)
    и \(\mathsf{level}(c') > \mathsf{level}(c)\).
\end{itemize}

\section*{Отчёт по криптостойкости профилей}

Для CKKS-профиля вычисляется суммарный размер коэффициентного модуля:
\[
    \log_2(Q)
    \approx
    \sum_i \mathsf{coeff\_modulus\_bits}_i.
\]

Для каждого профиля библиотека сравнивает это значение с лимитами
Microsoft SEAL:
\[
    \log_2(Q) \leq
    \mathsf{CoeffModulus::MaxBitCount}(N,\mathsf{sec\_level}).
\]

В отчёте фиксируются проверки для уровней:
\[
    \mathsf{tc128},\quad \mathsf{tc192},\quad \mathsf{tc256}.
\]

Для текущих профилей:
\[
\begin{array}{|c|c|c|c|c|c|}
\hline
\text{Профиль} & N & \log_2(Q) & \mathsf{tc128} & \mathsf{tc192} & \mathsf{tc256} \\
\hline
\texttt{basic\_ckks} & 8192 & 200 & \mathsf{PASS} & \mathsf{FAIL} & \mathsf{FAIL} \\
\hline
\texttt{depth\_ckks} & 16384 & 280 & \mathsf{PASS} & \mathsf{PASS} & \mathsf{FAIL} \\
\hline
\end{array}
\]

Общий уровень проекта определяется минимальным уровнем среди используемых
профилей:
\[
    \min(\mathsf{tc128}, \mathsf{tc192}) = \mathsf{tc128}.
\]

\section*{Демонстрационные приложения}

В проекте поддерживаются следующие запускные сценарии:

\[
\begin{array}{|l|p{9cm}|}
\hline
\text{Приложение} & \text{Назначение} \\
\hline
\texttt{demo\_secure\_stats} & Защищённое вычисление суммы и среднего. \\
\hline
\texttt{demo\_galois\_key\_optimization} & Сравнение полного и ограниченного набора Galois-ключей. \\
\hline
\texttt{demo\_abft} & Проверка ABFT-инвариантов для add, sub, mul и rotate. \\
\hline
\texttt{demo\_noise\_growth} & Анализ расхода глубины и роста ошибки. \\
\hline
\texttt{demo\_bootstrap\_pipeline} & Отчёт bootstrapping-модуля. \\
\hline
\texttt{bench\_ckks} & Измерение времени, ошибки и размеров ciphertext/ключей. \\
\hline
\texttt{bench\_chain\_accuracy} & Калибровка влияния chain length, scale и рабочей битности на точность. \\
\hline
\texttt{bench\_bootstrap\_parts} & Измерение plaintext-умножения, линейного преобразования, суммы слотов и полинома. \\
\hline
\texttt{bench\_parallel\_throughput} & Измерение параллельной обработки независимых ciphertext. \\
\hline
\texttt{demo\_profile\_report} & Таблица используемых CKKS-профилей. \\
\hline
\texttt{demo\_security\_report} & Проверка профилей по лимитам Microsoft SEAL. \\
\hline
\end{array}
\]

\end{document}
